\subsection{Usage in Experiments and Reproducibility Notes}
\label{app:prompts_usage}

All 120 prompts were executed against all five candidate models (\cref{sec:candidate_models}) using the standardized adapter interface (\cref{subsubsec:adapter_interface}) with fixed temperature ($T=0.3$; \cref{subsec:temperature_sweep}), system prompt template (strict Socratic; \cref{subsec:few_shot_prompts}), and context injection strategy (ephemeral user-message; \cref{subsec:contextualization_test}). Prompt order was randomized per model with a fixed seed to mitigate temporal confounds (transient provider load, diurnal latency variation) while preserving reproducibility (\cref{subsec:determinism}). The exact randomization seed and per-model prompt sequences are documented in the harness execution logs stored in \texttt{data/} (\cref{subsubsec:data_persistence}).

For context-bearing prompts (Adaptive Contextualization category), the \texttt{teacher\_context} field was prepended to the user message with a standardized delimiter (``Context: \{context\}\textbackslash{}n\textbackslash{}n'') to ensure syntactic consistency across models. Templating was performed via Jinja2 with no variable substitution in this study; template variables were reserved for future extensibility to parameterized problem generators (\cref{subsec:datasets_prompt_bank}).

Researchers seeking to replicate the benchmark should execute the harness with the configuration in \texttt{models.json} (per-model timeouts, rate limits, API endpoints; \cref{subsubsec:core_functionality}), environment variables in \texttt{.env.example} (API keys, logging levels), and prompt files in \texttt{data/test\_prompts.json} and \texttt{data/system\_prompts.json} (\cref{subsec:datasets_prompt_bank}). The complete Python environment is specified in \texttt{requirements.txt} (pinned dependency versions; \cref{subsec:determinism}). Validation checks at harness startup (\cref{subsubsec:config_validation}) surface configuration mismatches; diagnostics are logged to \texttt{logs/validation.log}.

Human rating procedures, including rater training protocols, blinding strategies, and adjudication thresholds, are documented in \cref{subsec:sampling_blinding}. The full Socratic Quality Rubric and Context Fidelity Rubric with scoring anchors and examples appear in \cref{app:socratic_quality_rubric,app:context_fidelity_rubric}. Raw model outputs, rubric scores, automated metrics, and aggregated results are available in CSV and SQLite formats in \texttt{data/} (\cref{subsubsec:data_persistence}); schema definitions appear in \cref{app:csv_schema,app:sqlite_schema}.
